\section{问题重述}

\subsection{问题背景}
在自主移动机器人导航中，单一传感器定位存在精度不足、鲁棒性差等问题。实际应用采用多源异构传感器组合方案，通过融合不同原理、不同频率的定位数据提升精度与实时性。本题提供两类定位数据——方式1（4\,Hz）与方式2（5\,Hz），存在起始时间不同步、采样频率不同、数据含噪声等问题，需融合输出 10\,Hz 高精度轨迹并完成沿途任务。

\subsection{问题提出}
\begin{itemize}
    \item \textbf{问题一}：附件1两类数据无噪声但存在设备开机时间先后不同，建立时间对齐模型，给出时间偏差及 10\,Hz 位置轨迹。
    \item \textbf{问题二}：附件2两类数据存在随机噪声和固定系统偏差，建立数据对齐与融合模型，给出时间偏差、系统偏差及 10\,Hz 位置轨迹。
    \item \textbf{问题三}：附件3为实际测量数据，判断是否存在系统偏差，并在此基础上完成问题二的任务。
    \item \textbf{问题四}：机器人沿附件3轨迹运动，对沿途目标执行射击或拍照任务，优化设计任务执行方案，尽可能多地完成任务。
\end{itemize}
